\subsection{SoC 设计}

本设计基于 ChipLab 提供的龙芯实验板 SoC。
Vivado 2023.2 工程位于 \texttt{fpga/loongson/2023.2}，顶层为 \texttt{soc\_top}；用户处理器以 \texttt{core\_top} 模块接入。
SoC 主要包含 DDR3、SPI Flash、UART、NAND Flash、Ethernet MAC、CONFREG（LED/数码管/按键/计时器）以及调试 UART 等。
CPU 经 AXI4 主口，通过时钟域转换与交叉开关访问各从设备。

\figplaceholder{fig:soc}{SoC 设计架构图}

\subsubsection{地址分配}

SoC 对各外设的物理地址映射如表~\ref{tab:addr} 所示（解码逻辑见 ChipLab AXI slave mux）。

\begin{table}[H]
\centering
\caption{地址映射表}
\label{tab:addr}
\small
\begin{tabular}{llll}
\toprule
地址段 & 起始地址 & 结束地址 & 说明 \\
\midrule
DDR3 Memory & \texttt{0x00000000} & \texttt{0x0FFFFFFF} & 主存（有效容量以 MIG 为准） \\
SPI / Boot 别名 & \texttt{0x1C000000} & \texttt{0x1C0FFFFF} & 启动相关窗口 \\
Configuration Registers & \texttt{0x1FD00000} & \texttt{0x1FD0FFFF} & GPIO / 板级计时器等 \\
UART Controller & \texttt{0x1FE00000} & \texttt{0x1FE0FFFF} & UART16550（数据口 \texttt{0x1FE001E0}） \\
NAND Controller & \texttt{0x1FE70000} & \texttt{0x1FE7FFFF} & NAND Flash \\
SPI Controller & \texttt{0x1FE80000} & \texttt{0x1FE8FFFF} & SPI Flash 控制 \\
Ethernet Controller & \texttt{0x1FF00000} & \texttt{0x1FF0FFFF} & 以太网 MAC \\
\bottomrule
\end{tabular}
\end{table}

软件常用虚地址窗口下，串口控制台对应 \texttt{0xBFE001E0}，内核链接/入口多为 \texttt{0xA0000000}，TFTP 临时加载缓冲使用 \texttt{0xA3000000}。

\subsubsection{外设控制器}

外设控制器主要沿用 ChipLab 提供的 IP，中断接入 CPU 的 \texttt{intrpt} 低位，配置见表~\ref{tab:periph}。

\begin{table}[H]
\centering
\caption{控制器配置及来源}
\label{tab:periph}
\begin{tabular}{lll}
\toprule
外设 & 控制器来源 & 中断位（示意） \\
\midrule
Ethernet & ChipLab & HWI0 \\
UART & ChipLab & HWI1 \\
SPI Flash & ChipLab & HWI2 \\
NAND Flash & ChipLab & HWI3 \\
DMA & ChipLab & HWI4 \\
CONFREG & ChipLab & --- \\
\bottomrule
\end{tabular}
\end{table}
